\documentclass[11pt]{article}
\usepackage[margin=1in]{geometry}
\usepackage{amsmath,amssymb}
\setlength{\parindent}{0pt}
\setlength{\parskip}{0.5em}

\begin{document}
\begin{center}
    \textbf{Compositional Difficulty}
\end{center}

\paragraph{1. Task and machine.}
Let $\Sigma$ be an alphabet. A task instance is
\[
\tau=(x,Y_x),\qquad x\in\Sigma^*,\qquad Y_x\subseteq\Sigma^*,
\]
where $x$ is the minimal immutable problem specification and $Y_x$
is the family of acceptable answers. Correctness is relational:
\[
\operatorname{correct}(y,x)\iff y\in Y_x.
\]
The machine is parameterized by $\theta$ and produces answers according to
$M_\theta(y\mid x,c)$, where $c$ contains all auxiliary context.
The correspondence with a Turing machine is: $x$ is the read-only input
tape, $c$ is the work tape, $\theta$ is the program, and $Y_x$ specifies
accepting outputs.

\paragraph{2. Context transformations.}
Initially $c_0$ is empty unless auxiliary information is supplied.
An append-only context update has the form
\[
c_i=c_{i-1}\oplus z_i,
\]
where $\oplus$ denotes concatenation. The update $z_i$ may contain an
intermediate deduction, retrieved fact, tool result, partial calculation,
subgoal, previous answer, or observation. Retrieval context and
chain-of-thought steps therefore have the same formal status: both are
information written into $c$.

Fix a class $\Phi$ of admissible primitive transformations
\[
\phi:(x,c)\mapsto(x,c'),\qquad c'=\phi(x,c).
\]
A depth-$d$ computation is a trace
\[
\pi:\quad
c_0\xrightarrow{\phi_1}c_1
\xrightarrow{\phi_2}\cdots
\xrightarrow{\phi_d}c_d,
\qquad \phi_i\in\Phi,
\]
followed by an answer sampled or decoded from
$M_\theta(\cdot\mid x,c_d)$.

A machine realization uses a context-generating transition $U_\theta$
and a halting/output operation $H_\theta$:
\[
c_{t+1}=c_t\oplus U_\theta(x,c_t),
\qquad
y=H_\theta(x,c_d).
\]
The computation is repeated composition of the context update.

\paragraph{3. Compositional depth.}
For error tolerance $\epsilon$, define
\[
D_\epsilon(x;\theta,\Phi)
=
\min\left\{
d:
\begin{array}{l}
\exists\,\phi_1,\ldots,\phi_d\in\Phi,\quad
c_i=\phi_i(x,c_{i-1}),\\
M_\theta(Y_x\mid x,c_d)\geq 1-\epsilon
\end{array}
\right\}.
\]
Here $M_\theta(Y_x\mid x,c_d)$ is the total probability of an acceptable
answer. Depth is the minimum number of primitive compositions needed
to make an answer recoverable with the specified reliability.

At the task level, let $T:x\mapsto Y_x$ and fix a primitive task class
$\mathcal P$. A realization is a composition of smaller tasks:
\[
T=T_d\circ\cdots\circ T_1,\qquad T_i\in\mathcal P.
\]
Its minimum depth is
\[
D_{\mathcal P}(T)
=
\min\left\{
d:T=T_d\circ\cdots\circ T_1,\quad T_i\in\mathcal P
\right\}.
\]
This is depth relative to a basis, as in circuit complexity. If arbitrary
functions are allowed as primitives, every task can be assigned depth one.
The primitive class must therefore be fixed before comparing depths.

Two scales should be distinguished:
\[
\begin{aligned}
d_{\mathrm{machine}}
&=\text{number of primitive machine transitions},\\
d_{\mathrm{comp}}
&=\text{number of dependent subtasks}.
\end{aligned}
\]
A subtask may require many machine transitions. For an autoregressive
model, machine transitions might correspond to token-generation steps;
compositional depth instead counts stages such as
$x\xrightarrow{T_1}z_1\xrightarrow{T_2}z_2\xrightarrow{T_3}y$.

\paragraph{4. Work space.}
For a trace $\pi$, the context-space usage is
\[
S(\pi)=\max_{0\leq i\leq d}|c_i|.
\]
Let $\Pi_\epsilon(x;\theta,\Phi)$ denote the admissible traces satisfying
the success condition above. Minimum work-space requirement is
\[
S_\epsilon(x;\theta,\Phi)
=
\min_{\pi\in\Pi_\epsilon(x;\theta,\Phi)}
\max_{0\leq i\leq|\pi|}|c_i|.
\]
If irrelevant context can be discarded, the corresponding live-space
measure is
\[
S_{\epsilon,\mathrm{live}}(x;\theta,\Phi)
=
\min_{\pi\in\Pi_\epsilon(x;\theta,\Phi)}
\max_i|\operatorname{live}(c_i)|,
\]
where $\operatorname{live}(c_i)$ is the information that must remain
available for later steps. Depth and space measure, respectively,
sequential composition and maintained working context.

\paragraph{5. Retrieval and prior examples.}
Retrieval is a context update $c_{i+1}=c_i\oplus r_i$.
If the origin of an update matters, label each transition by
\[
\kappa_i\in
\{\text{internal},\text{retrieval},\text{tool},\text{observation}\}.
\]
The resulting typed trace allows quantities such as
\[
D_R(\pi)=\#\{i:\kappa_i=\text{retrieval}\}
\]
without changing the underlying computation model.

Previously seen examples $X$ enter in either of two places.
Training examples are summarized by
$\theta=\operatorname{Learn}(X)$.
Examples available during inference are included in $c_0$.
Thus prior experience is represented through $\theta$ or $c$.

\paragraph{6. Noise and robust depth.}
Let perturbation channels act on the input and work context:
\[
\widetilde{x}\sim N_{x,\eta}(x),
\qquad
\widetilde{c}_i\sim N_{c,\eta}(c_i),
\]
where $\eta$ specifies the noise magnitude or distribution.
Define robust depth by applying the same success criterion to
computations using the perturbed input and context:
\[
D_{\epsilon,\eta}(x;\theta,\Phi)
=
\min\left\{
d:
\begin{array}{l}
\exists\text{ an admissible depth-}d\text{ computation},\\
\Pr_{\theta,\eta}(\widehat y\in Y_x)\geq1-\epsilon
\end{array}
\right\}.
\]
Probability is over the perturbations and machine output; correctness
is measured against the original answer family $Y_x$.
Nominal depth is $D_{\epsilon,0}$, and
\[
D_{\epsilon,\eta}-D_{\epsilon,0}
\]
can be interpreted as fragility under the specified noise.

\paragraph{7. Answer-family uncertainty.}
Having $|Y_x|>1$ need not make a task difficult: several answers may
be valid. Input perturbations can instead change which answers are valid,
from $Y_x$ to $Y_{\widetilde{x}}$. For a distance $\operatorname{dist}$
between answer sets, define
\[
U_\eta(x)
=
\mathbb E_{\widetilde{x}\sim N_{x,\eta}(x)}
\left[
\operatorname{dist}(Y_x,Y_{\widetilde{x}})
\right].
\]
This measures sensitivity of the answer relation to perturbations,
separately from computational depth and space.

\paragraph{8. Interpretation.}
Instance difficulty is the minimum depth and space of a successful
context-generating computation, relative to $\theta$ and a fixed
primitive basis $\Phi$; robustness describes its behavior under noise.
Reasoning chains contribute composition depth. Interacting constraints
can increase space and depth. Distractors can increase search, and
lexical complexity can require transformations before useful state is
constructed. Rare knowledge may require a retrieval transition.
Context length bounds the live information available but does not by
itself determine difficulty. Ambiguity concerns the answer relation
and its sensitivity rather than an additional computational resource.

\paragraph{9. Task families and composition.}
A task sampler requires a sampler for each component family, a rule for
composing compatible tasks, and a distribution over compositions.
Let
\[
\mathcal F=\{\mathcal F_1,\ldots,\mathcal F_m\},
\qquad
T\sim\mu_j(\,\cdot\mid k),
\]
where $\mu_j$ samples tasks from family $\mathcal F_j$ with structural
parameter $k$. Depending on the family, $k$ may specify the number of
constraints, dependency depth, or amount of state.

Represent a component task as a relation whose signature may depend on $k$:
\[
T:A_j^1(k)\times\cdots\times A_j^{m_j}(k)\longrightarrow 2^{B_j(k)},
\]
where $T(u)$ is the set of acceptable outputs for the input tuple $u$.
Every $T$ in the support of $\mu_j(\cdot\mid k)$ is total:
$T(u)\neq\emptyset$ for every input tuple $u$. A family whose natural
relation is partial restricts its input types to the inputs where it is
defined.

Components are composed along a directed acyclic graph (DAG) $G$. Each
node $v$ is either a leaf, which holds an input value of some type, or a
component node, which holds a task $T_v$ drawn from family $j_v$ with
parameter $k_v$ and has type $B_{j_v}(k_v)$. The $i$-th input of a
component node is an edge from another node whose type is contained in
$A_{j_v}^i(k_v)$, and the edges form no cycle. A value may be read by
several nodes. One node, the sink, holds the answer. An assignment gives
every node $v$ a value $z_v$ such that each component node satisfies
\[
z_v\in T_v(z_{v_1},\ldots,z_{v_m}),
\]
where $v_1,\ldots,v_m$ are its inputs. With the leaf values fixed, the
acceptable answers are the sink values of all assignments, and totality
guarantees that at least one assignment exists. A chain is the special
case in which every component has one input and every value has one
reader. There the definition reduces to
\[
(T_2\circ T_1)(u)
=
\bigcup_{z\in T_1(u)}T_2(z):
\]
an output is acceptable if it can be reached through an acceptable
intermediate output.

\paragraph{10. Sampling an instance.}
The sampler takes an answer type $C$, a budget $\beta$, a live-space cap
$\sigma$, a map $\lambda:C\to\Lambda$ that sorts answers into finitely many
classes, and a distribution $\omega$ on $\Lambda$. Each pair $(j,k)$ has a
positive integer cost $\delta_{jk}$, which bounds the depth in the
primitive basis $\mathcal P$ of every task that $\mu_j(\cdot\mid k)$
produces, and a weight $w_{jk}>0$. Each type $A$ that may appear at a leaf
has an input distribution $\nu_A$ and a weight $w_A>0$.

The sampler generates the structure backward from the answer and computes
the values forward from the inputs. Both halves condition a process on
how it ends.

The structure is built one node at a time, from the sink toward the
leaves, so every node is placed after all the nodes that read it. The
state is the pending set $Q$: the multiset of required types of the
values that placed nodes will read but whose producers have not been
placed. Initially $Q=\{C\}$, the answer. Each step takes one of two
actions:
\begin{itemize}
\item close a pending value of required type $A$ with a leaf: $Q$ loses
$A$; cost $0$, weight $w_A$;
\item produce a pending value of required type $A$ with a component node
for a pair $(j,k)$ with $B_j(k)\subseteq A$: $Q$ loses $A$, and each
input $i$ either opens a new pending value of required type $A_j^i(k)$
or attaches to a value already in $Q$, whose required type becomes its
intersection with $A_j^i(k)$; cost $\delta_{jk}$, weight $w_{jk}$.
\end{itemize}
An input that attaches to a pending value shares that value with its
other readers. Choosing different pending values counts as different
actions. A step may not make $|Q|$ exceed $\sigma$. The structure is
complete when $Q$ is empty and the budget is spent. Every node is placed
to produce a value that an already placed node reads, so every node feeds
the answer.

Reversing the order of placement gives an evaluation order. The pending
set at each point of the construction is the set of values live at the
same point of the evaluation: computed, and still to be read. So
$\sigma$ caps the live context of paragraph 4, measured in values. The
same structures could be built forward, with the live set as the state,
using the same states and an equivalent table. Building from the answer
starts at the one fixed type, $C$, and makes the tree case of paragraph
11 a grammar.

The structure must end complete. This condition involves only the
pending set and the remaining budget, so a table enforces it exactly.
The weight of a sequence of actions is the product of their weights, and
the total weight of the ways to complete the structure from state
$(Q,\beta')$ is
\[
Z(Q,\beta')
=
[\,Q=\emptyset,\ \beta'=0\,]
+
\sum_{a}w_a\,Z\big(Q_a,\beta'-\delta_a\big),
\]
where $a$ ranges over the actions allowed in state $Q$ with
$\delta_a\leq\beta'$, $Q_a$ is the resulting pending set, and $[\cdot]$
is $1$ when its condition holds and $0$ otherwise. Leaves cost nothing
and shrink $Q$, and required types range over intersections of input
types, a finite set; with the cap $\sigma$, the table has one entry per
multiset of at most $\sigma$ required types and remaining budget.
Restricting the allowed pairs $(j,k)$ or leaf types removes the
corresponding actions.

The values must end in an answer class $\ell$. Enforcing this exactly
would require a table indexed by intermediate values, so the sampler
repeats forward runs until one ends in $\ell$. Values run forward because
answers are computed from inputs; each family therefore needs only its
forward sampler. The answer class is controlled
because components that map many inputs to one output concentrate the
answers as the DAG grows, which makes guessing the most frequent answer
easier for larger tasks. Controlling the class of the answer suffices to
prevent this. The sampler proceeds as follows.
\begin{enumerate}
\item Sample the structure. Starting from $Q=\{C\}$ with budget
$\beta$, choose each action with probability
\[
\Pr\big(a\mid Q,\beta'\big)
=
\frac{w_a\,Z\big(Q_a,\beta'-\delta_a\big)}{Z(Q,\beta')}
\]
until the structure is complete. The result $s$ is a DAG with an
evaluation order whose component costs sum to $\beta$ and whose live set
never exceeds $\sigma$ values. When $Z(\{C\},\beta)>0$, every such
structure is drawn with probability proportional to the product of its
action weights, and no state along the way is a dead end.
\item Sample an answer class $\ell\sim\omega$.
\item Sample $T_v\sim\mu_{j_v}(\cdot\mid k_v)$ independently for every
component node and $z_v\sim\nu_A$ for every leaf of type $A$, and compute
the remaining values in evaluation order,
$z_v\in T_v(z_{v_1},\ldots,z_{v_m})$, where each $z_v$ is an acceptable
output produced by the family. Repeat this step until
$\lambda(z_\star)=\ell$, where $z_\star$ is the sink value.
\item Construct the instance:
$x=\operatorname{Encode}\big(s,(T_v)_v,(z_v)_{v\ \mathrm{leaf}}\big)$, and
$Y_x$ is the set of acceptable answers defined in paragraph 9.
\end{enumerate}
Totality keeps the forward run from getting stuck, and $z_\star\in Y_x$.
Write $\rho$ for the run, consisting of all $T_v$ and $z_v$,
$p_{\mathrm{fwd}}(\rho\mid s)$ for the distribution of a single run of
step 3 with structure $s$, and $p_s$ for the distribution of
$\lambda(z_\star)$ under $p_{\mathrm{fwd}}$. Steps 2 and 3 draw
\[
p(\rho\mid s)
=
p_{\mathrm{fwd}}(\rho\mid s)\,
\frac{\omega(\lambda(z_\star))}{p_s(\lambda(z_\star))},
\]
which is forward sampling with the distribution of answer classes replaced
by $\omega$. Within each class, runs follow the forward distribution
exactly. The expected number of forward runs is
$\sum_\ell\omega(\ell)/p_s(\ell)$; it is large when forward runs with
structure $s$ rarely reach a class that $\omega$ asks for.

The structure includes its evaluation order, which determines the live
sets. A DAG with several evaluation orders is therefore drawn once per
order, and $\operatorname{Encode}$ can present the nodes in that order, as
a program presents its statements. Here $\operatorname{Encode}$ produces a
problem statement containing the information needed to specify the task.
Generating labeled examples also requires a way to produce acceptable
outputs or check them.

The computation from step 3 provides intermediate results that can be
recorded in $c$. Some or all of this context may be supplied to the
solver, or the solver may be required to construct it.

\paragraph{11. Recursive composition.}
A component that is itself a DAG can be inlined into the outer DAG, so
recursive composition stays within DAGs. Two special cases recover simpler
forms. When every component has one input, the DAG is a chain. When every
value has exactly one reader, no input attaches to a pending value, the
DAG is a tree rooted at the answer, and the structure sampler is a
probabilistic grammar: a request for a value of type $A$ is replaced by a
component and requests for its inputs, which share the remaining budget.
Attaching is the one action the DAG sampler adds to the grammar.
Grammar-based generation has precedent
in compositional benchmarks: COGS generates examples using a probabilistic
context-free grammar.\footnote{Najoung Kim and Tal Linzen.
\emph{COGS: A Compositional Generalization Challenge Based on Semantic
Interpretation}. Proceedings of EMNLP, 2020, pp.\ 9087--9105.}
Values with several readers are what trees cannot express. They let live
space vary separately from the budget: the same $\beta$ can be spent on a
chain that keeps one value live or on a wide DAG that keeps many.

\paragraph{12. Construction structure and solving difficulty.}
The sampler controls the structure used to generate a task. This structure
need not be the shortest way to solve it. For example, for an invertible
function $f$,
\[
T=f^{-1}\circ f
\]
contains two components but reduces to the identity.

Let $T$ be the task computed by the DAG. Each component node has a
realization of depth at most $\delta_{j_vk_v}$ in the fixed primitive
basis $\mathcal P$, so evaluating the nodes in the sampled order
establishes
\[
D_{\mathcal P}(T)\leq\sum_{v}\delta_{j_vk_v}=\beta.
\]
The budget is an upper bound on depth. Equality requires an argument
excluding shortcuts. Components that are difficult individually may
simplify when composed. The structure bounds the other quantities of
paragraphs 3 and 4 in the same way. The number of component nodes on the
longest path from a leaf to the sink bounds the number of dependent
subtasks $d_{\mathrm{comp}}$, and the largest live set along the sampled
order, at most $\sigma$, bounds the number of values that must be held
between components.

The sampler therefore returns the instance together with its construction:
\[
\boxed{
\mathsf{Sample}(C,\beta,\sigma,\lambda,\omega)
\longrightarrow (x,Y_x,\gamma),
}
\]
where $\gamma$ records the structure $s$ and the component tasks. This
permits sampling under a controlled budget, live space, and distribution
of answer classes while leaving the relationship between construction
structure and minimum solving difficulty to be established.

\end{document}